\documentclass[10pt,a4paper]{amsart}
\usepackage[margin=19mm]{geometry}
\usepackage{amsmath,amssymb,mathtools}
\usepackage[scr=rsfs,cal=euler]{mathalfa}
\usepackage{xfrac}
\usepackage[unicode,hidelinks,bookmarks=false]{hyperref}
\newcommand{\ie}{i.e.\ }
\newcommand{\cf}{cf.\ }
\newcommand{\resp}{respectively\ }
\newcommand{\Subsection}{\subsection}
\providecommand{\nobreakdash}{\nobreak\hbox{-}}
\setlength{\emergencystretch}{2em}
\multlinegap0pt
\usepackage{amssymb}
\usepackage{mathtools}
\usepackage[shortlabels]{enumitem}
\usepackage{xcolor}
\newtheorem{theorem}{Theorem}
\newcommand{\CC}{\mathbb C}
\title{\bfseries Fock-Space Geometry and Demailly's Conjecture}
\author{Trung Hoa Dinh}
\date{Source: arXiv:2609.05866v2 (CC BY 4.0)}
\begin{document}
\maketitle

\noindent\emph{Source and licence.} \bfseries Fock-Space Geometry and Demailly's Conjecture. Version arXiv:2609.05866v2 (CC BY 4.0), distributed by arXiv under the Creative Commons Attribution 4.0 International licence, http://creativecommons.org/licenses/by/4.0/, as recorded in the arXiv licence metadata for this exact version and shown on the abstract page https://arxiv.org/abs/2609.05866v2. Full licence notice: \texttt{licenses/arxiv-2609.05866-CC-BY-4.0.txt}.\par\smallskip

\noindent\textit{Editorial scope.} Counted unit: the article's theorem thm:amplification (source0.tex lines 308--325). The article's complete original proof of it is delivered with it (source0.tex lines 327--409, one \texttt{proof} environment of 83 lines). No mathematics is written, repaired or re-derived: the statement and the proof are the article's own lines, in the article's own notation, and no retained line is substituted or re-flowed.  Retained uncounted as the source-local context the proof needs: lines 300--303.  Nothing was omitted from any retained span, so the package carries no omission record.  No retained line contains a citation, so the wrapper carries no bibliography and no bibliography entry.  No retained line contains a figure environment, an image-inclusion command or drawing code, so no image is delivered and the package declares no asset dependency.  Editorial formatting only, in addition: an AMS \texttt{amsart} preamble that restates the article's own declarations and macros used by the retained lines, namely the environments \texttt{theorem} on separate counters, together with the standard packages those lines need.  The retained environments print in this document as Theorem 1 in order of appearance, whereas the article prints the same items with the numbering of its own full document.  The complete native source of the article as distributed in the arXiv e-print for this exact version is bundled unchanged as \texttt{sources/209468/source0.tex}.
\begin{equation}
  S_A=\sum_i L_i^{r_i}S_{u_i-1},
  \label{eq:sourcecoverage}
\end{equation}

\begin{theorem}[Bosonic no-dark-state amplification]
\label{thm:amplification}
There exists a prime $p$ such that, for every $q=p^e$ and every positive multiplicity profile $\mathbf v=(v_1,\ldots,v_s)$ satisfying
\[
  v_i\ge q u_i+(n-1)(q-1)
  \qquad\text{for all }i,
\]
one has
\begin{equation}
  \alpha\!\left(\bigcap_iI(P_i)^{v_i}\right)
  \ge q\,\alpha(F_{\mathbf u})+(n-1)(q-1).
  \label{eq:amplification}
\end{equation}
Equivalently, if the source constraint operator has no dark state at particle number $A=a-1$, then the amplified target constraints have no dark state at
\[
  D_q=qA+n(q-1).
\]
\end{theorem}

\begin{proof}
Equation~\eqref{eq:sourcecoverage} is the surjectivity of the finite multiplication map
\[
  \Phi_A:\bigoplus_i S_{u_i-1}\longrightarrow S_A,
  \qquad
  (h_i)_i\longmapsto\sum_iL_i^{r_i}h_i.
\]
Choose a maximal minor witnessing surjectivity.  Spread the finitely many coefficients occurring in this matrix over a finitely generated integral $\mathbb Z$-algebra and localize so that the chosen determinant remains invertible.  A standard finite-type specialization argument then gives a finite residue field $\kappa$ of some characteristic $p>0$ for which the reduced source map remains surjective.  This prime is fixed from now on.

Let $q=p^e$ and put
\[
  D_q=qA+n(q-1).
\]
Over $\kappa$ set
\[
  J=(L_1^{r_1},\ldots,L_s^{r_s}).
\]
The source coverage says $J_A=S_A$, hence $J_t=S_t$ for every $t\ge A$.  We claim
\begin{equation}
  (J^{[q]})_{D_q}=S_{D_q},
  \qquad
  J^{[q]}=(L_1^{qr_1},\ldots,L_s^{qr_s}).
  \label{eq:frobcoverage}
\end{equation}
Because $\kappa$ is finite and hence perfect, every $G\in S_{D_q}$ has a unique decomposition
\begin{equation}
  G=\sum_{0\le\lambda_j<q} h_\lambda^{\,q}y^\lambda.
  \label{eq:frobdecomp}
\end{equation}
For each nonzero homogeneous summand,
\[
  q\deg h_\lambda=D_q-|\lambda|.
\]
There are $n+1$ variables and $|\lambda|\le(n+1)(q-1)$, so
\[
  q\deg h_\lambda
  \ge qA+n(q-1)-(n+1)(q-1)
  =q(A-1)+1.
\]
Since $\deg h_\lambda$ is integral, $\deg h_\lambda\ge A$.  Hence $h_\lambda\in J$, so $h_\lambda^q\in J^{[q]}$.  Every summand of \eqref{eq:frobdecomp} therefore lies in $J^{[q]}$, proving \eqref{eq:frobcoverage}.

In occupation-number language, the same step is the componentwise decomposition
\[
  \alpha=q\beta+\lambda,
  \qquad 0\le\lambda_j<q.
\]
The block $q\beta$ scales, while $\lambda$ is a bounded occupation residue.  The worst possible residue has total size $(n+1)(q-1)$; the shift in $D_q$ is exactly what prevents this residue from pushing the scalable core below the original full-coverage degree.  This is an arithmetic decomposition in a positive-characteristic Bargmann model, not a physical cloning or quantum channel.

The equality \eqref{eq:frobcoverage} is again a finite surjectivity statement.  Some maximal minor of its multiplication matrix is nonzero over $\kappa$; therefore the corresponding integral determinant is nonzero before reduction and remains nonzero after returning to characteristic zero.  Thus the amplified coverage also holds over $\CC$.

Now put
\[
  G_{\mathbf v}=\bigcap_iI(P_i)^{v_i}.
\]
If $v_i>D_q$ for some $i$, then $(G_{\mathbf v})_{D_q}=0$ immediately.  Otherwise define
\[
  e_i=D_q-v_i+1.
\]
The assumed lower bound on $v_i$ gives
\[
  e_i\le q(a-u_i)=qr_i.
\]
Hence
\[
  L_i^{qr_i}S_{D_q-qr_i}
  \subseteq
  L_i^{e_i}S_{D_q-e_i}.
\]
Summing and using the amplified coverage yields
\[
  S_{D_q}=\sum_iL_i^{e_i}S_{D_q-e_i}.
\]
Apolarity rotates this back to
\[
  (G_{\mathbf v})_{D_q}=0.
\]
Therefore
\[
  \alpha(G_{\mathbf v})\ge D_q+1
  =qa+(n-1)(q-1),
\]
which is \eqref{eq:amplification}.
\end{proof}

\end{document}
